\section{Secant objects on the split member}\label{sec:secantobjects}

\subsection{What the object must be, and what it cannot be}
\label{ssec:whatobject}

\Cref{thm:semiregclosure} asks for one perfect complex at one base point.
The cheapest candidate is a sum of line bundles on a split member, because
there the Weil classes are already polynomials in $\beta$, $\widehat\beta$ and
$\ell$ by \Cref{thm:splitclosed}, so the Chern character one wants is
manifestly in reach. This subsection settles that candidate. It does not
exist, and the obstruction is not semiregularity: it is the Chern character
condition, and it fails because the norm form of $K$ is positive
definite.\footnote{The norm form has now closed three different doors in this
paper, by three different mechanisms, and they should be kept apart.
\Cref{thm:character} uses the norm \emph{character} of $\Kx$ and an
equivariance hypothesis. \Cref{prop:divfails} uses the fact that the part $R$
of $H^{2}$ on which $\Kx$ acts through $\tau^{2}$ and $\bbar\tau^{\,2}$ has no
Hodge classes at a general point, which is a statement about the Hodge group.
\Cref{thm:positivity} below uses neither: it uses that $\Nm_{K/\QQ}$ is a
positive definite quadratic form on $K$ and that the multiplicities of a
filtration are positive integers. The three are independent, and only the
third is available at a special point, which is exactly where the object of
\Cref{prop:ktheory} has to live.}

\begin{lemma}[The semiregularity map of a split object]\label{lem:splitsemireg}
Let $Y$ be a smooth projective complex variety and let
\[
  E=\bigoplus_{i=1}^{s}L_{i}[p_{i}]^{\oplus m_{i}},
  \qquad \lambda_{i}=c_{1}(L_{i}),
\]
with the $L_{i}$ line bundles, the $\lambda_{i}$ pairwise distinct and
$p_{i}\in\ZZ$. Then for $\zeta=(\zeta_{ij})\in\Ext^{2}(E,E)$,
\begin{equation}\label{eq:splitsigma}
  \sigma_{E}(\zeta)_{q}
  =\frac{1}{q!}\sum_{i=1}^{s}(-1)^{p_{i}}\,
    \mathrm{tr}(\zeta_{ii})\,\lambda_{i}^{\,q},
\end{equation}
so $\sigma_{E}$ annihilates $\Ext^{2}(L_{i}[p_{i}],L_{j}[p_{j}])$ for $i\ne j$
and the traceless part of each diagonal block. Consequently $\sigma_{E}$ is
injective if and only if all three of the following hold:
\begin{enumerate}[label=\textup{(\roman*)},nosep]
\item $m_{i}=1$ for every $i$;
\item $H^{\,2+p_{i}-p_{j}}\bigl(Y,L_{j}L_{i}^{-1}\bigr)=0$ for all $i\ne j$;
\item the map
  \[
    \Phi\colon H^{2}(\cO_{Y})^{\oplus s}\longrightarrow H^{*}(Y,\CC),
    \qquad
    (\zeta_{i})\longmapsto\sum_{i}e^{\lambda_{i}}\zeta_{i},
  \]
  is injective.
\end{enumerate}
\end{lemma}

\begin{proof}
The Atiyah class of a direct sum is the direct sum of the Atiyah classes, and
that of a line bundle is its first Chern class, so $\mathrm{At}(E)$ is
diagonal with entries $\lambda_{i}$ and $\mathrm{At}(E)^{q}$ is diagonal with
entries $\lambda_{i}^{\,q}$. The matrix product $\zeta\circ\mathrm{At}(E)^{q}$
therefore has $(i,j)$ entry $\zeta_{ij}\lambda_{i}^{\,q}$, and the trace keeps
only the diagonal; a shift by $p$ multiplies the trace by $(-1)^{p}$. That is
\eqref{eq:splitsigma}, and it exhibits the two spaces the map kills.

For the equivalence, the off-diagonal part of $\Ext^{2}(E,E)$ is
$\bigoplus_{i\ne j}\Ext^{2}(L_{i}[p_{i}],L_{j}[p_{j}])
=\bigoplus_{i\ne j}H^{2+p_{i}-p_{j}}(L_{j}L_{i}^{-1})$, which must vanish,
giving (ii); the diagonal block at $i$ is $H^{2}(\cO_{Y})\otimes
\mathrm{End}(\CC^{m_{i}})$ and only its trace survives, so $m_{i}=1$, giving
(i). What is left is $\bigoplus_{i}H^{2}(\cO_{Y})$, and by
\eqref{eq:splitsigma} an element $(\zeta_{i})$ of it dies exactly when
$\sum_{i}\lambda_{i}^{\,q}\zeta_{i}=0$ for every $q\ge0$, which is (iii) after
absorbing the signs $(-1)^{p_{i}}$ into the $\zeta_{i}$ and the factorials
into the grading.
\end{proof}

Condition (ii) is cheap, because the shifts $p_{i}$ are at our disposal: only
their parities are fixed by the Chern character, and a line bundle whose first
Chern class is nondegenerate of index $k$ has cohomology in degree $k$ alone,
so one may choose the $p_{i}$ to avoid every index that occurs. Condition
(iii) is a rank. Condition (i) is free. What is not free is the Chern
character itself, and that is what the rest of this subsection is about.

\begin{notation}[Coordinates on the divisor classes]\label{not:splitcoords}
Let $A_{0}=X\times\widehat{X}$ be a split member, $n\ge2$, and keep
$\gamma=d\beta-\widehat{\beta}$ from \Cref{not:threeclasses}. Write
\[
  \theta^{\pm}=\gamma\mp\sqrt{-d}\,\ell\in\textstyle\bigwedge^{2}V_{\pm},
  \qquad
  R_{0}=\QQ\gamma\oplus\QQ\ell\subset R .
\]
Every $\theta\in R_{0}$ is $\theta=u\theta^{+}+\bbar u\,\theta^{-}$ for a
unique $u\in K$, and in coordinates $u=g+h\sqrt{-d}$ corresponds to
$\theta=2g\gamma+2dh\,\ell$. By \Cref{prop:splitgeom}(iv) the group
$\NS(A_{0})_{\QQ}$ at a split member with $X$ general is
$\QQ\eta\oplus R_{0}$, so every divisor class is $e\eta+\theta$ with
$e\in\QQ$ and $\theta\in R_{0}$, and we write $u(\lambda)\in K$ for the
coordinate of its $R_{0}$-part.
\end{notation}

\begin{lemma}[A basis in degree four]\label{lem:sixmonomials}
For $n\ge2$ the six products
$\beta^{2},\widehat\beta^{2},\ell^{2},\beta\widehat\beta,\beta\ell,
\widehat\beta\ell$
are linearly independent in $H^{4}(A_{0},\QQ)$. Equivalently the six classes
$\eta^{2}$, $\eta\theta^{+}$, $\eta\theta^{-}$, $(\theta^{+})^{2}$,
$(\theta^{-})^{2}$, $\theta^{+}\theta^{-}$ are a basis of the degree four part
of the subring generated by the divisor classes.
\end{lemma}

\begin{proof}
Grade $H^{*}(A_{0},\QQ)=\bigwedge^{*}U\otimes\bigwedge^{*}U^{\vee}$ by
bidegree. Then $\beta$ has bidegree $(2,0)$, $\widehat\beta$ has $(0,2)$ and
$\ell$ has $(1,1)$, so the six products have bidegrees $(4,0)$, $(0,4)$,
$(2,2)$, $(2,2)$, $(3,1)$, $(1,3)$. All are distinct except for the pair
$\beta\widehat\beta$ and $\ell^{2}$, and those two are independent in
$\bigwedge^{2}U\otimes\bigwedge^{2}U^{\vee}$ because
$\beta\widehat\beta=\beta\otimes\widehat\beta$ is decomposable there while
$\ell^{2}=-2\sum_{i<j}(x_{i}x_{j})\otimes(\xi_{i}\xi_{j})$ has all
$\binom{2n}{2}$ terms and $\binom{2n}{2}>n$ for $n\ge2$. The second
description follows because the change of basis between the two sets of six
classes is invertible over $K$: $\eta=d\beta+\widehat\beta$ and
$\gamma=d\beta-\widehat\beta$ span the same plane as
$\beta$ and $\widehat\beta$ since $\det\begin{psmallmatrix}d&1\\
d&-1\end{psmallmatrix}=-2d\ne0$, and $\theta^{\pm}$ span the same plane
as $\gamma$ and $\ell$ over $\CC$.
\end{proof}

The next lemma says the same thing at every abelian variety of Weil type, not
only at a split member, and it says it without coordinates. It is the fact the
positivity argument really turns on.

\begin{lemma}[$\theta^{+}\theta^{-}$ is never proportional to $\eta^{2}$]
\label{lem:tensorrank}
Let $(A,\iota,\eta)$ be of $(K,-1,n)$-Weil type with $n\ge2$, and let
$0\ne\theta^{\pm}\in\bigwedge^{2}V_{\pm}$. Then $\theta^{+}\theta^{-}$ is not a
scalar multiple of $\eta^{2}$ in $H^{4}(A,\CC)$.
\end{lemma}

\begin{proof}
Both classes lie in the summand
$\bigwedge^{2}V_{+}\otimes\bigwedge^{2}V_{-}$ of $H^{4}(A,\CC)$, the isotypic
component for the character $\Nm^{2}$, so the question is one of tensor rank
there. The polarisation class $\eta$ lies in $V_{+}\otimes V_{-}$ and is the
hermitian form of \Cref{prop:weilline}, hence nondegenerate: in dual bases
$\eta=\sum_{a=1}^{2n}v_{a}\wedge w_{a}$ with $v_{a}$ a basis of $V_{+}$ and
$w_{a}$ a basis of $V_{-}$. Squaring,
\[
  \eta^{2}
  =\sum_{a,b}(v_{a}\wedge w_{a})(v_{b}\wedge w_{b})
  =-2\sum_{a<b}(v_{a}v_{b})\otimes(w_{a}w_{b}),
\]
which is diagonal in the induced bases of $\bigwedge^{2}V_{+}$ and
$\bigwedge^{2}V_{-}$ and therefore has tensor rank exactly $\binom{2n}{2}$.
The class $\theta^{+}\theta^{-}$ is, up to sign,
$\theta^{+}\otimes\theta^{-}$, of tensor rank one. For $n\ge2$ one has
$\binom{2n}{2}\ge6>1$, so the two cannot be proportional.
\end{proof}

\begin{theorem}[The positivity obstruction]\label{thm:positivity}
Let $(A,\iota,\eta)$ be of $(K,-1,n)$-Weil type with $n\ge2$, and suppose that
the only divisor classes on which $\Kx$ acts through the norm character are
the multiples of $\eta$, that is
\begin{equation}\label{eq:noextranorm}
  \NS(A)_{\QQ}\cap\bigl(V_{+}\otimes V_{-}\bigr)=\QQ\,\eta .
\end{equation}
Let $E$ be a perfect complex on $A$ whose Chern character is effective in the
exponential basis, meaning
\begin{equation}\label{eq:effectivech}
  \ch(E)=\sum_{i=1}^{s}m_{i}\,e^{\lambda_{i}},
  \qquad m_{i}\in\QQ_{>0},\quad\lambda_{i}\in\NS(A)_{\QQ},
\end{equation}
and write $\lambda_{i}=e_{i}\eta+\theta_{i}$ with $\theta_{i}$ in the part $R$
of $H^{2}(A,\QQ)$ carrying the characters $\tau^{2}$ and $\bbar\tau^{\,2}$, and
$u_{i}\in K$ for the $\theta^{+}$-coordinate of $\theta_{i}$. If $\ch_{2}(E)$
is a flat Hodge class, then
\begin{equation}\label{eq:normsum}
  \sum_{i=1}^{s}m_{i}\,\Nm_{K/\QQ}(u_{i})=0 ,
\end{equation}
and since $\Nm_{K/\QQ}$ is positive definite and every $m_{i}$ is positive,
$u_{i}=0$ for every $i$. Hence every $\lambda_{i}$ is a rational multiple of
$\eta$, and $\ch_{k}(E)\in\QQ\,\eta^{k}$ for every $k$. In particular
$\ch_{n}(E)$ has no component on the Weil line.
\end{theorem}

\begin{proof}
From \eqref{eq:effectivech},
$\ch_{2}(E)=\tfrac12\sum_{i}m_{i}\lambda_{i}^{2}$. Decompose $H^{4}(A,\CC)$
into $\Kx$-isotypic pieces and keep the one for $\Nm^{2}$, namely
$\bigwedge^{2}V_{+}\otimes\bigwedge^{2}V_{-}$. Writing
$\theta_{i}=u_{i}\theta^{+}+\bbar u_{i}\theta^{-}$ for a fixed choice of
generators $\theta^{\pm}$ of $\bigwedge^{2}V_{\pm}$, one has
\[
  \lambda_{i}^{2}
  =e_{i}^{2}\eta^{2}
   +2e_{i}\eta\theta_{i}
   +u_{i}^{2}(\theta^{+})^{2}+\bbar u_{i}^{2}(\theta^{-})^{2}
   +2\Nm_{K/\QQ}(u_{i})\,\theta^{+}\theta^{-},
\]
and of these five terms only the first and the last are in the $\Nm^{2}$
piece: $\eta\theta_{i}$ carries $\Nm\tau^{2}$ and $\Nm\bbar\tau^{\,2}$, and
$(\theta^{\pm})^{2}$ carry $\tau^{4}$ and $\bbar\tau^{\,4}$, all four distinct
from $\Nm^{2}$. Hypothesis \eqref{eq:noextranorm} is what makes the
decomposition $\lambda_{i}=e_{i}\eta+\theta_{i}$ available with $e_{i}$ a
scalar. So the $\Nm^{2}$-part of $\ch_{2}(E)$ is
\[
  \tfrac12\Bigl(\sum_{i}m_{i}e_{i}^{2}\Bigr)\eta^{2}
  +\Bigl(\sum_{i}m_{i}\Nm_{K/\QQ}(u_{i})\Bigr)\theta^{+}\theta^{-}.
\]

Now compare with the flat Hodge classes of degree four. For $n\ge3$ they are
$\QQ\eta^{2}$ by \Cref{thm:hdgcount}. For $n=2$ they are
$\QQ\eta^{2}\oplus\HW(A)$, and $\HW(A)$ spans the summands for $\tau^{4}$ and
$\bbar\tau^{\,4}$, which are not $\Nm^{2}$. In both cases the $\Nm^{2}$-part
of a flat Hodge class of degree four is a multiple of $\eta^{2}$. By
\Cref{lem:tensorrank} the class $\theta^{+}\theta^{-}$ is not such a multiple,
so its coefficient must vanish, which is \eqref{eq:normsum}.

Finally $\Nm_{K/\QQ}(g+h\sqrt{-d})=g^{2}+dh^{2}$ with $d>0$, so
$\Nm_{K/\QQ}(u)\ge0$ with equality only at $u=0$, and a sum of nonnegative
terms with positive weights vanishes only when every term does. Then
$\ch_{k}(E)=\frac{1}{k!}\bigl(\sum_{i}m_{i}e_{i}^{k}\bigr)\eta^{k}$, and
$\eta^{k}$ is not on the Weil line by \Cref{thm:hdgcount}.
\end{proof}

\begin{corollary}[The cheapest objects do not exist]\label{cor:nofiltered}
Let $A$ satisfy \eqref{eq:noextranorm}, which holds at every split member of
\Cref{thm:splitclosed} and at every quaternionic base point of
\Cref{thm:quatweil}.\footnote{At a split member $\NS(A)_{\QQ}$ has rank three
with basis $\beta,\widehat\beta,\ell$ by \Cref{prop:splitgeom}(iv), and of
these $\eta$ spans the norm-character part while $\gamma$ and $\ell$ span
$R$. At a quaternionic point $\NS(A)_{\QQ}$ is the space of Rosati-symmetric
elements of $B=(-d,b)_{\QQ}$, which is spanned by $1$, $j$ and $ij$; the first
gives $\eta$ and the other two are semilinear over $K$, hence lie in $R$. In
both cases the norm-character part of $\NS(A)_{\QQ}$ is exactly $\QQ\eta$.}
Then \Cref{thm:semiregclosure} cannot be applied at $A$ with any of the
following:
\begin{enumerate}[label=\textup{(\alph*)},nosep]
\item a coherent sheaf admitting a filtration by line bundles, in particular
  any direct sum of line bundles;
\item a simple semi-homogeneous vector bundle in the sense of Mukai
  \cite{Muk78}, since such a bundle of rank $r$ has $\ch(E)=r\,e^{c_{1}(E)/r}$;
\item any iterated extension of the objects in (a) and (b).
\end{enumerate}
In each case the failure is not of the semiregularity hypothesis but of the
Chern character hypothesis of \Cref{prop:ktheory}: no such $E$ has a Chern
character which is a flat section of Hodge classes with a nonzero Weil
component.
\end{corollary}

\begin{proof}
Each of (a), (b) and (c) has a Chern character of the shape
\eqref{eq:effectivech} with positive multiplicities, by additivity of $\ch$
along a filtration and by Mukai's computation of the Chern character of a
simple semi-homogeneous bundle. Apply \Cref{thm:positivity}.
\end{proof}

\begin{remark}[Where the theorem has force, and why it does not contradict
Markman]\label{rem:markmanpositivity}
At a very general member of the family $\NS(A)_{\QQ}=\QQ\eta$, so
\Cref{thm:positivity} is true for the empty reason that there are no
$\theta_{i}$ to begin with. Its content is at the special members, which is
exactly where \Cref{thm:semiregclosure} needs its object to live and where the
base points of \Cref{thm:everyfamily} are. That is the opposite of the usual
situation in this paper, where the special members are the easy ones.

The theorem also does not touch the objects of Markman \cite{Mar25a,Mar25b}.
Those are sheaves whose class in $K_{0}$ is not a positive combination of
exponentials of divisor classes, so \eqref{eq:effectivech} fails for them at
the first step; and they live at very general points, where as just noted the
statement is vacuous anyway. What the theorem does is remove the class of
objects one would reach for first, and it removes it for a reason which is
visible before any cohomology is computed.
\end{remark}

\begin{proof}
Immediate from \Cref{thm:positivity}, since the conclusion there contradicts
the requirement that $\ch_{n}(E)$ have a nonzero component on the Weil line.
\end{proof}

\begin{remark}[What is left of the candidate, and what it must look like]
\label{rem:whatitmustbe}
\Cref{thm:positivity} uses positivity of the multiplicities and nothing else,
so it says exactly where to go next: the object must be a complex whose
class in $K_{0}$ is not effective, that is one in which the line bundles
appear with multiplicities of both signs. Such objects exist. Writing
$E=F\oplus G[1]$ with $F$ and $G$ sums of line bundles, the class in $K_{0}$
is $[F]-[G]$ and \eqref{eq:normsum} becomes
$\sum_{L_{i}\subset F}\Nm(u_{i})=\sum_{L_{j}\subset G}\Nm(u_{j})$, which is
satisfiable; and \Cref{lem:splitsemireg} still applies verbatim, with the
shifts absorbing condition (ii).

Two further facts, both computational, say how little room is left at $n=2$,
and how much more there is beyond it. First, the map $\Phi$ of
\Cref{lem:splitsemireg}(iii) is injective at $n=2$ for $s\le3$ and, for a
general choice of $\lambda_{1},\dots,\lambda_{4}$, has rank $22$ rather than
$24$; so at $n=2$ no more than three line bundles may occur at all. This is
special to $n=2$: at $n=3$ the same map is injective for a general choice of
up to ten classes, and the dimension count permits up to $33$. Second, an
exhaustive search over the classes $a\beta+b\widehat\beta+c\ell$ with
$|a|,|b|,|c|\le3$, for $d=1$ and $d=2$, finds no signed tuple of at most
three line bundles whose Chern character is a flat section of Hodge classes
with a nonzero Weil component; the tuples of at most four that do satisfy the
Chern character condition with a nonzero Weil component, $888$ of them at
$d=1$ and $30$ at $d=2$, all have exactly four terms, and all of them fail
\Cref{lem:splitsemireg}(iii): the kernel has dimension $2$, $3$, $4$ or $6$
at $d=1$ and $3$ or $6$ at $d=2$, never $0$. Item (XIII) of
\Cref{sec:verification} carries both computations.

So the object of \Cref{thm:semiregclosure}, if it exists at a split member, is
not a sum of line bundles in any arrangement of signs and shifts: it must have
$\Ext^{2}$ that the trace does not see, that is, it must be indecomposable in
the derived category. That is a sharper description of the remaining
construction than \Cref{rem:tworemain} gives, and it is as far as we have been
able to push it.
\end{remark}

\subsection{The object, written down, and what it costs}
\label{ssec:pte}

\Cref{prop:ktheory} produced the object of \Cref{thm:semiregclosure} by an
existence argument: the Chern character is an isomorphism onto the Chow ring
with rational coefficients, so a $K$-theory class with the right character
exists. That is enough for the theorem and tells one nothing about the object.
This subsection writes it down. The answer turns the question into a
Diophantine one, and the Diophantine problem has a classical name.

\begin{setupenv}[Objects supported on $R$]\label{setup:pte}
Let $A$ be of $(K,-1,n)$-Weil type, $n\ge2$, satisfying
\eqref{eq:noextranorm}, and keep $\theta^{\pm}$ from
\Cref{not:splitcoords}. For $u\in K$ with $\theta(u)=u\theta^{+}+\bbar
u\,\theta^{-}\in\NS(A)_{\QQ}$, and for multiplicities $m_{i}\in\ZZ$, consider
\[
  v=\sum_{i=1}^{s}m_{i}\,e^{\theta(u_{i})},
  \qquad
  T_{q,r}=\sum_{i=1}^{s}m_{i}\,u_{i}^{q}\,\bbar u_{i}^{\,r}
  \quad (q,r\ge0).
\]
\end{setupenv}

\begin{proposition}[The Chern character in terms of the $u_{i}$]
\label{prop:chintermsofu}
In \Cref{setup:pte},
\begin{equation}\label{eq:chexpansion}
  v_{k}
  =\sum_{\substack{q+r=k\\ 0\le q,r\le n}}
    \frac{T_{q,r}}{q!\,r!}\,(\theta^{+})^{q}(\theta^{-})^{r},
\end{equation}
because $(\theta^{\pm})^{q}=0$ for $q>n$. Consequently $v$ is a flat section
of Hodge classes with a nonzero Weil component if
\begin{equation}\label{eq:pteconditions}
  T_{q,r}=0\ \ \text{for all } 0\le q,r\le n \text{ with }
  (q,r)\notin\{(0,0),(n,0),(0,n)\},
  \qquad T_{n,0}\ne0,
\end{equation}
and then $v_{0}=T_{0,0}=\sum_{i}m_{i}$, $v_{k}=0$ for $0<k<2n$ with $k\ne n$,
and
\[
  v_{n}=\frac{1}{n!}\Bigl(T_{n,0}(\theta^{+})^{n}
        +\bbar{T}_{n,0}(\theta^{-})^{n}\Bigr),
\]
a nonzero rational multiple of a generator of the Weil line, by
\Cref{thm:splitclosed}.
\end{proposition}

\begin{proof}
Expand $\theta(u_{i})^{k}$ by the binomial theorem, divide by $k!$ and collect;
the vanishing of $(\theta^{\pm})^{q}$ for $q>n$ is because
$\theta^{\pm}\in\bigwedge^{2}V_{\pm}$ and $\dim_{\CC}V_{\pm}=2n$. Under
\eqref{eq:pteconditions} the only surviving terms of \eqref{eq:chexpansion}
are the ones listed, and $(\theta^{+})^{n}$ spans $\bigwedge^{2n}V_{+}$, which
is the Weil line over $\CC$.
\end{proof}

So the search for the object of \Cref{thm:semiregclosure} is the search for a
finitely supported integer-valued function on $\cO_{K}$ that annihilates every
monomial $u^{q}\bbar u^{\,r}$ with $q,r\le n$ except three of them. Two things
follow at once, and they pull in opposite directions.

\begin{lemma}[Equal norms force the Prouhet--Tarry--Escott condition]
\label{lem:pte}
Suppose all the $u_{i}$ in \Cref{setup:pte} have the same norm
$N=\Nm_{K/\QQ}(u_{i})>0$. Then \eqref{eq:pteconditions} holds if and only if
\begin{equation}\label{eq:ptepure}
  \sum_{i}m_{i}u_{i}^{\,k}=0\ \ \text{for } k=0,1,\dots,n-1,
  \qquad
  \sum_{i}m_{i}u_{i}^{\,n}\ne0 .
\end{equation}
With $m_{i}=\pm1$ this says that the $u_{i}$ split into two multisets on the
circle $\Nm=N$ whose power sums agree through degree $n-1$ and differ at
degree $n$: a Prouhet--Tarry--Escott pair of degree $n-1$, constrained to a
norm circle of $\cO_{K}$.
\end{lemma}

\begin{proof}
For $q\ge r$ one has $u^{q}\bbar u^{\,r}=N^{r}u^{q-r}$, so
$T_{q,r}=N^{r}\sum_{i}m_{i}u_{i}^{\,q-r}$, and the pairs $(q,r)$ with
$q,r\le n$ realise every difference $k=q-r$ from $0$ to $n$. The excluded
pairs $(0,0)$ and $(n,0)$ are the ones with $r=0$ and $k=0$ or $k=n$; but
$k=0$ also occurs at $(1,1),\dots,(n,n)$, which are not excluded, so
$\sum_{i}m_{i}=0$ is required after all. The remaining differences
$k=1,\dots,n-1$ each occur at some non-excluded pair, and $k=n$ occurs only at
$(n,0)$.
\end{proof}

\begin{theorem}[Equal norms cannot be semiregular]\label{thm:equalnorms}
Let $E=\bigoplus_{i}L_{i}[p_{i}]$ realise \eqref{eq:ptepure} with all norms
equal. Then $\ch_{0}(E)=0$, and the semiregularity map $\sigma_{E}$ has a
kernel of dimension at least $n^{2}$.
\end{theorem}

\begin{proof}
By \Cref{lem:pte} the case $k=0$ of \eqref{eq:ptepure} gives
$\sum_{i}m_{i}=0$, which is $\ch_{0}(E)$. For the kernel, take
$\zeta_{i}=m_{i}\zeta$ for a single $\zeta\in H^{2}(\cO_{A})$; then by
\Cref{lem:splitsemireg}(iii)
\[
  \Phi\bigl((m_{i}\zeta)_{i}\bigr)
  =\Bigl(\sum_{i}m_{i}e^{\lambda_{i}}\Bigr)\zeta=\ch(E)\,\zeta ,
\]
which by \Cref{prop:chintermsofu} is a multiple of $\omega\,\zeta$ with
$\omega$ a generator of the Weil line. Now $\HW(A)_{\CC}$ is spanned by
$\bigwedge^{2n}V_{+}$ and $\bigwedge^{2n}V_{-}$, and the Weil type condition
says $\dim V_{\pm}^{0,1}=n$; writing
$H^{2}(\cO_{A})=\bigwedge^{2}(V_{+}^{0,1}\oplus V_{-}^{0,1})$, the product of
a generator of $\bigwedge^{2n}V_{+}$ with $\zeta$ kills the components of
$\zeta$ in $\bigwedge^{2}V_{+}^{0,1}$ and in
$V_{+}^{0,1}\otimes V_{-}^{0,1}$, and likewise with the roles exchanged. So
$\omega\,\zeta=0$ for every $\zeta$ in
$V_{+}^{0,1}\otimes V_{-}^{0,1}$, a space of dimension $n^{2}$, and all of
those lie in the kernel of $\Phi$.
\end{proof}

\begin{corollary}[At least $n+1$ norms]\label{cor:nplusone}
If $E$ as in \Cref{setup:pte} satisfies \eqref{eq:pteconditions} and
$\ch_{0}(E)\ne0$, then the $u_{i}$ take at least $n+1$ distinct norms.
\end{corollary}

\begin{proof}
Let $N_{1},\dots,N_{\nu}$ be the distinct norms and $M_{j}$ the sum of the
$m_{i}$ over the level $\Nm(u_{i})=N_{j}$. The conditions $T_{r,r}=0$ for
$r=1,\dots,n$ read $\sum_{j}M_{j}N_{j}^{\,r}=0$ for $r=1,\dots,n$. The matrix
$(N_{j}^{\,r})_{1\le r\le n,\,1\le j\le\nu}$ has rank $\min(n,\nu)$ because the
$N_{j}$ are distinct and nonzero, so $\nu\le n$ forces every $M_{j}=0$ and
hence $\ch_{0}(E)=\sum_{j}M_{j}=0$.
\end{proof}

\subsubsection*{The smallest solution, and what it does}

At $n=3$ and $K=\QQ(\sqrt{-3})$ the Prouhet--Tarry--Escott condition of
\Cref{lem:pte} has a solution of striking simplicity: the norm circle
$\Nm=1$ carries the six roots of unity of $K$, and the two cosets of the cube
roots of $-1$ and of $+1$ are a Prouhet--Tarry--Escott pair of degree
two.\footnote{Explicitly, with $\zeta$ a primitive sixth root of unity, the
cube roots of $-1$ are $\zeta,\zeta^{3},\zeta^{5}$ and those of $+1$ are
$1,\zeta^{2},\zeta^{4}$; each triple sums to zero and has zero sum of squares,
while the sums of cubes are $-3$ and $+3$. A search over the norm circles of
$\cO_{K}$ with $\Nm\le200$ and $|g|,|h|\le25$, for $K=\QQ(\sqrt{-d})$ with
$d=1,2,3,7$, finds this solution and its conjugate at $n=3$, one further pair
at $d=1$ and $n=4$, and nothing else; item (XIII) of \Cref{sec:verification}
records the search.}

\begin{proposition}[The object at $n=3$, written out]\label{prop:explicitobject}
Let $K=\QQ(\sqrt{-3})$, $n=3$, and let $A_{0}=X\times\widehat X$ be a split
member. In the basis $\beta,\widehat\beta,\ell$ of $\NS(A_{0})_{\QQ}$ put
\[
  \lambda_{1}=(3,-1,3),\quad
  \lambda_{2}=(-6,2,0),\quad
  \lambda_{3}=(3,-1,-3),
\]
and $\lambda_{3+i}=-\lambda_{i}$ for $i=1,2,3$, with $m_{i}=+1$ for
$i\le3$ and $m_{i}=-1$ for $i\ge4$. Then the six classes are distinct and
integral, the perfect complex
\[
  E=\bigl(L_{1}\oplus L_{2}\oplus L_{3}\bigr)\ \oplus\
    \bigl(L_{4}\oplus L_{5}\oplus L_{6}\bigr)[3]
\]
has
\[
  \ch_{k}(E)=0\ \ (k\ne3),
  \qquad
  \ch_{3}(E)=12\,\omega_{1},
\]
and every difference $\lambda_{j}-\lambda_{i}$ with $i\ne j$ is nondegenerate
of index exactly $3$, so that conditions (i) and (ii) of
\Cref{lem:splitsemireg} hold for this shift. Condition (iii) fails: the map
$\Phi$ has rank $75$ instead of $90$, with kernel of dimension
$15=\dim H^{2}(\cO_{A_{0}})$.
\end{proposition}

\begin{proof}
Everything asserted is a finite exact computation, carried out in item (XIII)
of \Cref{sec:verification}. The classes are the images under
$\theta(u)=2g\gamma+2dh\,\ell$, $u=g+h\sqrt{-3}$, of the six roots of unity
of $K$ scaled by $2$; the Chern character is \Cref{prop:chintermsofu} with the
Prouhet--Tarry--Escott data above; the index computation and the rank of
$\Phi$ are the two remaining items. For the shift, condition (ii) asks that
$2+p_{i}-p_{j}\ne3$ whenever $i\ne j$, that is $p_{i}-p_{j}\ne1$; within each
of the two blocks the difference is $0$, and between them it is $\pm3$.
\end{proof}

\begin{remark}[What this settles and what it does not]\label{rem:pteoutcome}
Three things are settled. The object of \Cref{prop:ktheory} is not an
abstraction: at $n=3$ over $\QQ(\sqrt{-3})$ it is six line bundles and a
shift, and its Chern character is exactly $12$ times a Weil class with every
other graded piece zero. The construction is completely general, and
\Cref{lem:pte} identifies the constraint as a Prouhet--Tarry--Escott problem
on a norm circle, which is a question of a kind that can be attacked. And
\Cref{cor:nplusone} says what any semiregular solution must look like: the
$u_{i}$ must take at least $n+1$ distinct norms, because otherwise the rank is
zero and \Cref{thm:equalnorms} produces a kernel of dimension $n^{2}$.

The remaining question was the existence of a solution with $n+1$ norms. The
level sums are then forced up to scale, and they compete with the supply of
elements of $\cO_{K}$ at each level; \Cref{thm:pterange} carries out that
competition in a range, and \Cref{cor:noline} then shows that no solution in
any range could have served, because a line bundle summand off the
polarisation line is obstructed along the family whatever the norms.
\end{remark}

\subsubsection*{The obstruction classes of the explicit object}

\Cref{cor:kernelfamily} and \Cref{rem:whatmeasures} say that the kernel of
$\sigma_{E}$ contains a copy of the tangent space of the Weil family, and
that this is a statement about where the semiregularity map is blind, not a
statement that the object is obstructed. For the explicit object the
obstruction classes can be written down, and they are not zero. The
computation is elementary and we record it in general.

Recall the conventions of \Cref{ssec:tangent}: a first-order deformation of
$A$ is $v\in\Hom(H^{1,0},H^{0,1})=H^{1}(A,T_{A})$, extended by zero on
$H^{0,1}$ and acting on $H^{*}(A,\CC)$ as a derivation $v\lrcorner$; a class
$\xi$ of type $(p,p)$ stays of type $(p,p)$ to first order along $v$ exactly
when $v\lrcorner\,\xi=0$. For a line bundle $L$ with $\lambda=c_{1}(L)$ the
Atiyah class is $\lambda\in H^{1}(\Omega^{1}_{A})$, its contraction with $v$
is $v\lrcorner\,\lambda\in H^{2}(\cO_{A})=H^{0,2}(A)$, and this is the
obstruction to deforming $L$ along $v$; a shift does not change it. For
$E=\bigoplus_{i}L_{i}[p_{i}]$ the Atiyah class is diagonal with entries
$\lambda_{i}$, as in the proof of \Cref{lem:splitsemireg}, so the
obstruction to deforming $E$ along $v$ is the diagonal element
\begin{equation}\label{eq:obdef}
  \ob_{E}(v)=\bigl(v\lrcorner\,\lambda_{i}\bigr)_{i}
  \in\bigoplus_{i}H^{0,2}(A)\subset\Ext^{2}(E,E).
\end{equation}

\begin{proposition}[The obstruction map of a split object]
\label{prop:splitobstruction}
Let $A_{0}=X\times\widehat X$ be a split member with $\beta$ principal and
$X$ general, so that $\NS(A_{0})_{\QQ}=\langle\beta,\widehat\beta,\ell\rangle$
by \Cref{prop:splitgeom}(iv). Let $T=T_{[A_{0}]}\cD_{n,n}$ be the space of
polarised first-order deformations commuting with $K$, of dimension $n^{2}$
by \Cref{prop:firstorder}, and put
\[
  T_{\mathrm{spl}}
  =\bigl\{\,v\in\Hom(H^{1,0},H^{0,1})\;:\;
    v\lrcorner\,\beta=v\lrcorner\,\widehat\beta=v\lrcorner\,\ell=0\,\bigr\}.
\]
Let $E=\bigoplus_{i=1}^{s}L_{i}[p_{i}]$ with $\lambda_{i}=c_{1}(L_{i})$ and
let $\ob_{E}\colon T\to\bigoplus_{i}H^{0,2}(A_{0})$ be the map
\eqref{eq:obdef} restricted to $T$.
\begin{enumerate}[label=\textup{(\roman*)},nosep]
\item $T_{\mathrm{spl}}$ is the tangent space to the split locus: it has
  dimension $n(n+1)/2$, it is contained in $T$, and it consists of the
  deformations of $X\times\widehat X$ induced by the polarised deformations
  of $(X,\beta)$.
\item If the classes $\lambda_{i}$ together with $\eta$ span
  $\NS(A_{0})_{\QQ}$, then $\ker\ob_{E}=T_{\mathrm{spl}}$ and
  $\operatorname{rank}\ob_{E}=n(n-1)/2$, the codimension of the split locus
  in $\cD_{n,n}$.
\item If $\ch_{k}(E)\in\QQ\eta^{k}$ for $k\ne n$ and
  $\ch_{n}(E)\in\QQ\eta^{n}\oplus\HW(A_{0})$, then
  $\sigma_{E}\circ\ob_{E}=0$ on $T$: for every $v\in T$,
  \begin{equation}\label{eq:sigmaob}
    \sigma_{E}\bigl(\ob_{E}(v)\bigr)=v\lrcorner\,\ch(E)=0 .
  \end{equation}
\item For the object of \Cref{prop:explicitobject}, at $n=3$ and
  $K=\QQ(\sqrt{-3})$: $\ker\ob_{E}=T_{\mathrm{spl}}$ has dimension $6$ and
  $\ob_{E}$ has rank $3$; the image of $\ob_{E}$ lies in $\ker\sigma_{E}$;
  the subspace $\Delta=\{z\cdot\id_{E}:z\in\Ann_{H^{0,2}}(\omega_{1})\}$ of
  \Cref{cor:kernelfamily} has dimension $9$ and meets the image of $\ob_{E}$
  only in zero, so that $\Delta\oplus\image\ob_{E}$ is a twelve-dimensional
  subspace of the fifteen-dimensional kernel of $\sigma_{E}$; and
  $\ob_{E}(v)$ has components $v\lrcorner\,\lambda_{i}$ on $L_{i}$ and
  $-v\lrcorner\,\lambda_{i}$ on $L_{i}^{-1}[3]$, so it lies on the line of
  vectors of the shape $(z,\dots,z)$ only when it vanishes. In particular
  $E$ has a nonzero first-order obstruction along every $v\in T$ not tangent
  to the split locus, and no element $z\cdot\id_{E}$ of $\Ext^{2}(E,E)$
  cancels it.
\end{enumerate}
\end{proposition}

\begin{proof}
(i) Let $P=H^{1,0}(X)\subset H^{1}(X,\CC)=U^{\vee}_{\CC}$ with basis
$p_{1},\dots,p_{n}$, so that $p_{1},\dots,p_{n},\bbar p_{1},\dots,\bbar p_{n}$
is a basis of $U^{\vee}_{\CC}$, and let $q_{a},\bbar q_{a}$ be the dual basis
of $U_{\CC}=H^{1}(\widehat X,\CC)$; the dual of a conjugate basis is the
conjugate of the dual basis, so the notation is consistent. The Poincar\'e
class is the identity element of $U^{\vee}\otimes U$,
\[
  \ell=\sum_{a=1}^{n}\bigl(p_{a}\wedge q_{a}+\bbar p_{a}\wedge\bbar q_{a}\bigr),
\]
and it is of type $(1,1)$, which forces $H^{1,0}(\widehat X)=\langle\bbar
q_{a}\rangle$ and $H^{0,1}(\widehat X)=\langle q_{a}\rangle$: the Hodge
structure of $\widehat X$ is the dual one. Since $\beta$ is of type $(1,1)$
and nondegenerate, $\beta=\sum_{a,b}h_{ab}\,p_{a}\wedge\bbar p_{b}$ with
$h\in\GL_{n}(\CC)$; and since $\beta$ is principal with Gram matrix $B$ in a
symplectic basis of $U$ and $\widehat\beta$ has Gram matrix $B=-B^{-1}$ in
the dual basis of $U^{\vee}$, the Gram matrix of $\widehat\beta$ in the
basis $p_{a},\bbar p_{a}$ is minus the inverse of the Gram matrix
$\begin{psmallmatrix}0&h\\-h^{\top}&0\end{psmallmatrix}$ of $\beta$ in the
basis $q_{a},\bbar q_{a}$, which gives
$\widehat\beta=-\sum_{a,b}(h^{-1})_{ab}\,\bbar q_{a}\wedge q_{b}$.

Now $H^{1,0}(A_{0})=\langle p_{a}\rangle\oplus\langle\bbar q_{a}\rangle$ and
$H^{0,1}(A_{0})=\langle\bbar p_{a}\rangle\oplus\langle q_{a}\rangle$, so a
deformation $v$ is four $n\times n$ matrices,
\[
  v(p_{a})=\sum_{c}\bigl(A_{ca}\,\bbar p_{c}+C_{ca}\,q_{c}\bigr),
  \qquad
  v(\bbar q_{a})=\sum_{c}\bigl(B_{ca}\,\bbar p_{c}+D_{ca}\,q_{c}\bigr),
\]
and $v$ kills $\bbar p_{a}$ and $q_{a}$. Contracting,
\[
  v\lrcorner\,\ell
  =\sum_{a,c}\Bigl(A_{ca}\,\bbar p_{c}\wedge q_{a}+C_{ca}\,q_{c}\wedge q_{a}
   +B_{ca}\,\bbar p_{a}\wedge\bbar p_{c}+D_{ca}\,\bbar p_{a}\wedge q_{c}\Bigr),
\]
and reading off the three kinds of monomials, $v\lrcorner\,\ell=0$ if and only
if $D=-A^{\top}$ and $B$, $C$ are symmetric. Next
\[
  v\lrcorner\,\beta
  =\sum_{a,b,c}h_{ab}\Bigl(A_{ca}\,\bbar p_{c}\wedge\bbar p_{b}
   +C_{ca}\,q_{c}\wedge\bbar p_{b}\Bigr),
\]
so $v\lrcorner\,\beta=0$ if and only if $Ch=0$, that is $C=0$, and $Ah$ is
symmetric. Likewise $v\lrcorner\,\widehat\beta=0$ if and only if $B=0$ and
$Dh^{-1}$ is symmetric. Given $D=-A^{\top}$, the last condition reads
$A^{\top}h^{-1}=h^{-\top}A$, which after multiplication by $h^{\top}$ on the
left and by $h$ on the right is $h^{\top}A^{\top}=Ah$, the symmetry of $Ah$
again. Hence
\[
  T_{\mathrm{spl}}=\bigl\{(A,B,C,D)=(A,0,0,-A^{\top})\;:\;Ah\ \text{symmetric}\bigr\},
\]
of dimension $n(n+1)/2$. It is the tangent space to the split locus: the
same computation on $X$ alone shows that the polarised first-order
deformations of $(X,\beta)$ are the matrices $A$ with $Ah$ symmetric, the
space $T_{[X]}\cA_{n}$ of dimension $n(n+1)/2$; the deformation of
$\widehat X$ induced by $A$ is the dual one, $-A^{\top}$; and the map
$A\mapsto(A,0,0,-A^{\top})$ is injective. Finally $T_{\mathrm{spl}}\subseteq
T$: for $v\in T_{\mathrm{spl}}$, $v\lrcorner\,\eta=0$ because
$\eta\in\langle\beta,\widehat\beta\rangle$, so $v$ is polarised, and
$v\lrcorner\,\omega_{1}=v\lrcorner\,\omega_{2}=0$ because $\omega_{1}$ and
$\omega_{2}$ are polynomials in $\beta,\widehat\beta,\ell$ by
\Cref{thm:splitclosed} and $v\lrcorner$ is a derivation; by
\Cref{prop:firstorder} such a $v$ commutes with $K$.

(ii) Every $v\in T$ has $v\lrcorner\,\eta=0$. If the $\lambda_{i}$ and
$\eta$ span $\NS(A_{0})_{\QQ}=\langle\beta,\widehat\beta,\ell\rangle$, then
$v\lrcorner\,\lambda_{i}=0$ for all $i$ is equivalent, for $v\in T$, to
$v\lrcorner\,\beta=v\lrcorner\,\widehat\beta=v\lrcorner\,\ell=0$, so
$\ker\ob_{E}=T\cap T_{\mathrm{spl}}=T_{\mathrm{spl}}$ by (i), and the rank
is $n^{2}-n(n+1)/2=n(n-1)/2$, which is the codimension of
\Cref{prop:splitgeom}(iii).

(iii) By \eqref{eq:splitsigma}, $\sigma_{E}(\zeta)=\sum_{i}(-1)^{p_{i}}
\zeta_{ii}\,e^{\lambda_{i}}$ for a diagonal $\zeta$. Since $v\lrcorner$ is a
derivation and $v\lrcorner\,\lambda_{i}$ has even degree,
$v\lrcorner\,\lambda_{i}^{\,q}=q\,\lambda_{i}^{\,q-1}(v\lrcorner\,\lambda_{i})$
and so $v\lrcorner\,e^{\lambda_{i}}=(v\lrcorner\,\lambda_{i})\,e^{\lambda_{i}}$.
Therefore
\[
  \sigma_{E}\bigl(\ob_{E}(v)\bigr)
  =\sum_{i}(-1)^{p_{i}}(v\lrcorner\,\lambda_{i})\,e^{\lambda_{i}}
  =v\lrcorner\sum_{i}(-1)^{p_{i}}e^{\lambda_{i}}
  =v\lrcorner\,\ch(E).
\]
For $v\in T$ one has $v\lrcorner\,\eta^{k}=k\eta^{k-1}(v\lrcorner\,\eta)=0$,
and $v\lrcorner\,\omega=0$ for every $\omega\in\HW(A_{0})$ by
\Cref{prop:firstorder}; under the hypothesis on $\ch(E)$ this gives
\eqref{eq:sigmaob}.

(iv) The three classes $\lambda_{1}=\gamma+3\ell$, $\lambda_{2}=-2\gamma$,
$\lambda_{3}=\gamma-3\ell$ span $\langle\gamma,\ell\rangle$, and
$\langle\eta,\gamma\rangle=\langle\beta,\widehat\beta\rangle$ by the
determinant in the proof of \Cref{lem:sixmonomials}, so (ii) applies and
gives the kernel and the rank; \Cref{prop:explicitobject} gives the
hypothesis of (iii), so the image lies in $\ker\sigma_{E}$. The subspace
$\Delta$ has dimension $\dim\Ann_{H^{0,2}}(\omega_{1})=n^{2}=9$ by
\Cref{lem:annihilator}. The components of $\ob_{E}(v)$ on $L_{i}$ and on
$L_{i}^{-1}[3]$ are $v\lrcorner\,\lambda_{i}$ and
$v\lrcorner\,(-\lambda_{i})=-v\lrcorner\,\lambda_{i}$, so an element of the
image of the shape $(z,\dots,z)$ has $z=-z$, hence is zero; this is
$\Delta\cap\image\ob_{E}=0$, and with $\dim\ker\sigma_{E}=15$ from
\Cref{prop:explicitobject} the dimension count follows. Every number in (iv)
is also computed directly in item (XXV) of \Cref{sec:verification}, in the
model of item (XIII).
\end{proof}

\begin{remark}[What the obstruction classes say]\label{rem:obstructionmeaning}
\Cref{prop:splitobstruction} sharpens \Cref{rem:whatmeasures} in one
direction and closes a door in another. The direction: for the explicit
object the obstruction classes along the family are nonzero, and they are
nonzero exactly in the $n(n-1)/2$ directions normal to the split locus,
which are the directions in which the Weil class has to be transported.
The door: those classes lie in the kernel of $\sigma_{E}$ by
\eqref{eq:sigmaob}, and \eqref{eq:sigmaob} holds for every object whose
Chern character is a flat section of Hodge classes along the family, so the
semiregularity map cannot detect the obstruction of any such object. The
obstruction is therefore genuine, and it is invisible to the criterion. For
the explicit object the situation is the worst possible: the twelve
dimensions $\Delta\oplus\image\ob_{E}$ of the kernel are accounted for by
the two mechanisms of \Cref{cor:kernelfamily} and \eqref{eq:obdef}, the
remaining three lie on the even side of \Cref{cor:pteobstructed} as well,
and none of the fifteen is an artefact. The classes of the shape
$z\cdot\id_{E}$ are the ones a change of the twisting class would
contribute, one and the same $z$ to every rank one summand; the last
assertion of (iv) says that no such change removes the obstruction. So the
explicit object does not move off the split locus to first order, as a
complex and as a twisted complex, and the question of \Cref{q:residual} for
it is not one of verifying a hypothesis but of replacing the object.

The positive content is the recipe. By (ii), a split object whose classes
span $\NS(A_{0})_{\QQ}$ with $\eta$ has obstruction map of rank
$n(n-1)/2$ against the family, and the spanning hypothesis is more than is
needed: in the model of item (XXV) a single class $\theta\in\NS(A_{0})_{\QQ}$
not proportional to $\eta$ already has $\{v\in T:v\lrcorner\,\theta=0\}=
T_{\mathrm{spl}}$, for each of the seven classes tested at
$(n,d)=(2,1),(2,3),(3,3)$. So every split object built from line bundles
other than powers of $\eta$ is obstructed in exactly the normal directions
to the split locus, and \Cref{thm:positivity} says that an object built
from powers of $\eta$ carries no Weil class. What survives is the statement
of \Cref{rem:familyuniform}: the object must have positive rank.
\end{remark}

\subsubsection*{The evaluation map of the explicit object}

\Cref{thm:weakfails} places the copy $\Delta$ of the tangent space inside the
image of the evaluation map. For a split object the whole of
$\operatorname{ev}_{E}$ can be written down, and for the explicit object the
entire kernel of $\sigma_{E}$ turns out to lie in its image. We identify
$HH^{2}(A_{0})$ with
\begin{equation}\label{eq:HT2}
  HT^{2}(A_{0})
  =H^{0,2}(A_{0})\;\oplus\;\Hom\bigl(H^{1,0},H^{0,1}\bigr)\;\oplus\;
   \textstyle\bigwedge^{2}\bigl(H^{1,0}\bigr)^{\vee}
\end{equation}
by the Hochschild--Kostant--Rosenberg isomorphism, the three summands being
$H^{2}(\cO)$, $H^{1}(T)$ and $H^{0}(\bigwedge^{2}T)$ for the trivial tangent
bundle. An element $x=(z,v,\pi)$ acts on $H^{*}(A_{0},\CC)$ by $z\wedge$, by
the derivation $v\lrcorner$ of \Cref{ssec:tangent}, and by the double
contraction $\pi\lrcorner$, which removes two $(1,0)$ slots; we write
$x\lrcorner\,\xi$ for the sum of the three.%
\footnote{The double contraction is a second order operator. For a class
$\lambda$ of type $(1,1)$, which has a single $(1,0)$ slot,
$\pi\lrcorner\,\lambda=0$ and $\pi\lrcorner\,\lambda^{k}=\binom{k}{2}
\lambda^{k-2}\,(\pi\lrcorner\,\lambda^{2})$, because $\pi\lrcorner$ is a
composite of two interior products, each of which is an antiderivation, and
the only nonzero terms are those in which the two interior products hit two
different factors. Summing over $k$ with the factorials gives
$\pi\lrcorner\,e^{\lambda}=\tfrac12(\pi\lrcorner\,\lambda^{2})\,e^{\lambda}$,
exactly parallel to $v\lrcorner\,e^{\lambda}=(v\lrcorner\,\lambda)\,e^{\lambda}$
for the first order operator $v\lrcorner$. This is the whole content of the
compatibility \eqref{eq:sigmaev} below: the semiregularity map of a line
bundle is multiplication by $e^{\lambda}$, and every summand of $HT^{2}$ acts
on $e^{\lambda}$ by a scalar-valued factor which is precisely the value of
$\operatorname{ev}_{L}$. The same identity, for arbitrary $E$, is the
commutative square of Buchweitz and Flenner \cite[Corollary 4.3]{BF03} and
its Hochschild form in \cite[Lemma 11.3]{Mar25b}; on an abelian variety the
Todd class is trivial and no correction to the HKR map is needed
\cite{Cal05,BF08}.}

\begin{proposition}[The evaluation map of a split object]\label{prop:evaluation}
Let $A_{0}$ be a split member with $X$ general and let
$E=\bigoplus_{i=1}^{s}L_{i}[p_{i}]$ with $\lambda_{i}=c_{1}(L_{i})$ and all
$\Ext^{2}(L_{i}[p_{i}],L_{j}[p_{j}])$, $i\ne j$, zero, so that
$\Ext^{2}(E,E)=\bigoplus_{i}H^{0,2}(A_{0})$.
\begin{enumerate}[label=\textup{(\roman*)},nosep]
\item Under \eqref{eq:HT2}, $\operatorname{ev}_{E}$ is diagonal:
  \begin{equation}\label{eq:evsplit}
    \operatorname{ev}_{E}(z,v,\pi)_{i}
    =z+v\lrcorner\,\lambda_{i}+\tfrac12\,\pi\lrcorner\,\lambda_{i}^{2}
    \;\in\;H^{0,2}(A_{0})=\Ext^{2}\bigl(L_{i}[p_{i}],L_{i}[p_{i}]\bigr).
  \end{equation}
\item For every $x\in HT^{2}(A_{0})$,
  \begin{equation}\label{eq:sigmaev}
    \sigma_{E}\bigl(\operatorname{ev}_{E}(x)\bigr)=x\lrcorner\,\ch(E),
  \end{equation}
  and consequently
  $\ker\sigma_{E}\cap\image\operatorname{ev}_{E}
   =\operatorname{ev}_{E}\bigl(\{x:x\lrcorner\,\ch(E)=0\}\bigr)$.
  The weakened hypothesis of \cite[Question 11.4]{Mar25b} holds for $E$ if
  and only if $\operatorname{ev}_{E}$ kills every $x$ that preserves
  $\ch(E)$ to first order.
\item $\dim HT^{2}(A_{0})=2\binom{2n}{2}+4n^{2}$ and
  $\dim\Ext^{2}(E,E)=s\binom{2n}{2}$; for $s\ge5$ and every $n\ge2$ the
  second exceeds the first, so $\operatorname{ev}_{E}$ is not surjective and
  \cite[Lemma 11.3]{Mar25b} cannot apply to a sum of five or more line
  bundles.
\item For the object of \Cref{prop:explicitobject}, with
  $\lambda_{i}=u_{i}\theta^{+}+\bbar u_{i}\theta^{-}$ as in
  \Cref{not:splitcoords}: the space $\{x:x\lrcorner\,\ch(E)=0\}$ is
  $\Ann_{H^{0,2}}(\omega_{1})\oplus\{v:v\lrcorner\,\omega_{1}=0\}\oplus
  \{\pi:\pi\lrcorner\,\omega_{1}=0\}$, of dimensions $9$, $18$ and $9$; the
  first and the third summands are both carried by $\operatorname{ev}_{E}$
  onto $\Delta$; the second is carried onto
  \[
    \Sigma=\bigl\{\,(u_{i}a+\bbar u_{i}b)_{i=1}^{6}\;:\;
      a\in\textstyle\bigwedge^{2}V_{+}^{0,1},\ b\in\bigwedge^{2}V_{-}^{0,1}\,\bigr\},
    \qquad \dim\Sigma=n(n-1)=6,
  \]
  which contains the image of the obstruction map of
  \Cref{prop:splitobstruction}; $\Delta\cap\Sigma=0$; and
  \[
    \ker\sigma_{E}=\Delta\oplus\Sigma\subseteq\image\operatorname{ev}_{E}.
  \]
  So every direction in which $\sigma_{E}$ is blind is the value of
  $\operatorname{ev}_{E}$ on a Hochschild class that preserves $\ch(E)$ to
  first order, and the weakened hypothesis fails for $E$ by all fifteen
  dimensions. The rank of $\operatorname{ev}_{E}$ is $45$. The same
  containment $\ker\sigma_{E}\subseteq\image\operatorname{ev}_{E}$, with
  $\dim\ker\sigma_{E}=15$, holds for each of the eighteen configurations of
  \Cref{rem:familyuniform}.
\end{enumerate}
\end{proposition}

\begin{proof}
(i) The evaluation map is the action of $HH^{*}(A_{0})$ on
$\Ext^{*}(E,E)$, which under HKR is contraction of a polyvector in
$H^{q}(\bigwedge^{p}T)$ with $\tfrac{1}{p!}\mathrm{At}(E)^{p}$
\cite{BF08}; for $p=0$ this is $z\cdot\id_{E}$ as in the proof of
\Cref{thm:weakfails}, for $p=1$ it is $v\lrcorner\,\mathrm{At}(E)$, the map
\eqref{eq:obdef}, and for $p=2$ it is $\tfrac12\pi\lrcorner\,\mathrm{At}(E)^{2}$.
The Atiyah class of $E$ is diagonal with entries $\lambda_{i}$, so all three
are diagonal with the entries displayed.

(ii) By \eqref{eq:splitsigma}, $\sigma_{E}(\zeta)=\sum_{i}(-1)^{p_{i}}\zeta_{i}
e^{\lambda_{i}}$ for diagonal $\zeta$. The three factors in \eqref{eq:evsplit}
are exactly the scalars by which $z\wedge$, $v\lrcorner$ and $\pi\lrcorner$
act on $e^{\lambda_{i}}$: the first trivially, the second because
$v\lrcorner$ is a derivation and $v\lrcorner\,\lambda_{i}$ has even degree,
the third by the footnote above. Hence
$\sigma_{E}(\operatorname{ev}_{E}(x))=\sum_{i}(-1)^{p_{i}}\,x\lrcorner\,
e^{\lambda_{i}}=x\lrcorner\,\ch(E)$. The description of
$\ker\sigma_{E}\cap\image\operatorname{ev}_{E}$ follows: $\operatorname{ev}_{E}(x)$
lies in $\ker\sigma_{E}$ if and only if $x\lrcorner\,\ch(E)=0$.

(iii) The three summands of \eqref{eq:HT2} have dimensions $\binom{2n}{2}$,
$4n^{2}$ and $\binom{2n}{2}$. The inequality
$s\binom{2n}{2}>2\binom{2n}{2}+4n^{2}$ is $(s-2)(2n-1)>4n$, and
$4n/(2n-1)\le8/3$ for $n\ge2$, so it holds for $s\ge5$; it is checked for
$2\le n\le30$ and $5\le s\le40$ in \Cref{sec:lean}.

(iv) Since $\ch(E)=12\omega_{1}$ is of type $(3,3)$, the three summands of
$x\lrcorner\,\ch(E)$ have types $(3,5)$, $(2,4)$ and $(1,3)$, so the kernel
is the direct sum of the three kernels. The first is $\Ann_{H^{0,2}}(\omega_{1})$,
of dimension $n^{2}=9$ by \Cref{lem:annihilator}. The second is
$\Hom(V_{+}^{1,0},V_{+}^{0,1})\oplus\Hom(V_{-}^{1,0},V_{-}^{0,1})$, of
dimension $2n^{2}=18$, by the proof of \Cref{prop:firstorder}, whose
argument for (i)$\Leftrightarrow$(iii) does not use the polarisation. For the
third, decompose $\bigwedge^{2}(H^{1,0})^{\vee}$ as
$\bigwedge^{2}(V_{+}^{1,0})^{\vee}\oplus(V_{+}^{1,0})^{\vee}\otimes(V_{-}^{1,0})^{\vee}
\oplus\bigwedge^{2}(V_{-}^{1,0})^{\vee}$; here $\omega_{1}$ is a combination
of $(\theta^{+})^{3}$ and $(\theta^{-})^{3}$ with both coefficients nonzero,
$(\theta^{\pm})^{3}$ spans $\bigwedge^{3}V_{\pm}^{1,0}\otimes\bigwedge^{3}V_{\pm}^{0,1}$,
the outer blocks of $\pi$ act injectively on the respective factor
$\bigwedge^{3}V_{\pm}^{1,0}$ and land in different summands of $H^{1,3}$,
and the middle block kills both; so the third kernel is the middle block,
of dimension $n^{2}=9$.

At a split point $\theta^{\pm}$ are of type $(1,1)$, so $\theta^{\pm}\in
V_{\pm}^{1,0}\otimes V_{\pm}^{0,1}$, and they are nondegenerate pairings
because $(\theta^{\pm})^{3}\ne0$. Write $\lambda_{i}^{2}=u_{i}^{2}(\theta^{+})^{2}
+\bbar u_{i}^{2}(\theta^{-})^{2}+2\Nm(u_{i})\,\theta^{+}\theta^{-}$. A middle
block $\pi$ kills $(\theta^{\pm})^{2}$, which have no $V_{\mp}^{1,0}$ slot, so
$\tfrac12\pi\lrcorner\,\lambda_{i}^{2}=\Nm(u_{i})\,\pi\lrcorner\,(\theta^{+}\theta^{-})$;
the six $u_{i}$ of \Cref{prop:explicitobject} are $\pm\zeta_{6},\pm1,
\pm\zeta_{6}^{-1}$, all of norm one, so this value is independent of $i$ and
$\operatorname{ev}_{E}(\pi)=z_{\pi}\cdot\id_{E}$ with
$z_{\pi}=\pi\lrcorner\,(\theta^{+}\theta^{-})$. The map $\pi\mapsto z_{\pi}$
from the middle block to $V_{+}^{0,1}\otimes V_{-}^{0,1}$ is the tensor
product of the two nondegenerate pairings, hence an isomorphism onto the
middle block of $H^{0,2}$, which is $\Ann_{H^{0,2}}(\omega_{1})$ by
\Cref{lem:annihilator}. So $\operatorname{ev}_{E}$ carries the third kernel
onto $\Delta$, as it carries the first by \Cref{thm:weakfails}.

For $v=(v_{+},v_{-})$ in the second kernel,
$v\lrcorner\,\lambda_{i}=u_{i}\,v_{+}\lrcorner\,\theta^{+}+\bbar u_{i}\,
v_{-}\lrcorner\,\theta^{-}$ with $a=v_{+}\lrcorner\,\theta^{+}\in
\bigwedge^{2}V_{+}^{0,1}$ and $b=v_{-}\lrcorner\,\theta^{-}\in\bigwedge^{2}V_{-}^{0,1}$;
as $v_{\pm}$ ranges over $\Hom(V_{\pm}^{1,0},V_{\pm}^{0,1})$, $a$ and $b$
range over the whole of $\bigwedge^{2}V_{\pm}^{0,1}$ because the pairings
$\theta^{\pm}$ are nondegenerate, which is the surjectivity of
$\operatorname{alt}(Vc)$ in \eqref{eq:phidelta}. So the image is $\Sigma$.
The map $(a,b)\mapsto(u_{i}a+\bbar u_{i}b)_{i}$ is injective, because
$(u_{1},\bbar u_{1})$ and $(u_{2},\bbar u_{2})$ are linearly independent in
$\CC^{2}$, $u_{1}/\bbar u_{1}=\zeta_{6}^{2}$ differing from
$u_{2}/\bbar u_{2}=1$; hence $\dim\Sigma=2\binom{n}{2}=n(n-1)=6$. The
obstruction map of \Cref{prop:splitobstruction} is the restriction to the
polarised $K$-commuting $v$, so its image lies in $\Sigma$. An element of
$\Sigma$ of the shape $(z,\dots,z)$ has $u_{i}a+\bbar u_{i}b=z=-z$ on
comparing $i$ with $3+i$, so $z=0$ and then $a=b=0$: $\Delta\cap\Sigma=0$.
By (ii), $\Delta\oplus\Sigma\subseteq\ker\sigma_{E}\cap\image\operatorname{ev}_{E}$,
and its dimension $9+6=15$ equals $\dim\ker\sigma_{E}$ by
\Cref{prop:explicitobject}, so the three spaces coincide. The rank of
$\operatorname{ev}_{E}$ and the statement for the eighteen configurations
are the computations of item (XXVI) of \Cref{sec:verification}, which also
recomputes every dimension in this proof in the model of item (XIII).
\end{proof}

\begin{figure}[tbp]
\centering
  \includegraphics[width=\textwidth]{fig_evaluation}
\caption{\Cref{prop:evaluation} for the explicit object at $n=3$. Left: the
three summands of $HT^{2}(A_{0})$, drawn with heights proportional to their
dimensions, and inside each the part that preserves $\ch(E)$ to first order:
the annihilator of $\omega_{1}$, the deformations killing $\omega_{1}$ with
the polarised $K$-commuting tangent space $T$ in front, and the middle block
of bivectors. Right: $\Ext^{2}(E,E)$ as six stacked copies of $H^{0,2}$, one
per summand of $E$, and inside it the kernel of $\sigma_{E}$: the column
$\Delta$, one and the same class in every layer, and the column $\Sigma$,
whose entries depend on the summand through $u_{i}$. The three arrows are the
evaluation map on the three kernels; the two teal ones both land on
$\Delta$, the ochre one on $\Sigma$, and together they cover the kernel.
Below, the commutative square that makes this a statement about the weakened
criterion of \cite[Question 11.4]{Mar25b}.}
\label{fig:evaluation}
\end{figure}

\begin{remark}[Question 11.4 for split objects, and what is left of it]
\label{rem:question114}
For a split object the two criteria stand or fall together, and for the
explicit object both fall completely. \Cref{thm:weakfails} makes the first
$n^{2}$ dimensions of the failure uniform over every concentrated object in
every dimension; \Cref{prop:evaluation}(iv) accounts for the remaining
$n(n-1)$ at $n=3$ by the deformations of the two halves $V_{\pm}$ of the
Hodge structure separately, and the equality $\ker\sigma_{E}=\Delta\oplus\Sigma$
says that the two mechanisms exhaust the kernel: nothing in it is
accidental. Nothing here touches the case the question was asked for. A
sheaf of positive rank has $z\cdot\ch(E)\ne0$ for every $z\ne0$, so its
$\Delta$ is zero, and whether $\operatorname{ev}_{E}$ kills the deformations
that preserve $\kappa(E)$ is a statement about the sheaf and not about its
Chern character. Markman's candidate in his Section 12, a rank one complex on
a principally polarised abelian fourfold with Chern character on the secant
$\langle e^{\pm\sqrt{-d}\Theta}\rangle$, is exactly such a sheaf, and he
records that the verification remains to be done. \Cref{thm:factor} reduces
that verification to one identity on the fourfold, and
\Cref{thm:candidatefails} shows that the identity fails: the candidate does
not satisfy the weakened criterion either.
\end{remark}

\begin{theorem}[Where a semiregular split object can live]\label{thm:pterange}
Keep \Cref{setup:pte} and suppose $E$ satisfies \eqref{eq:pteconditions} with
$\ch_{0}(E)\ne0$, so that the $u_{i}$ take $\nu\ge n+1$ distinct norms by
\Cref{cor:nplusone}. Write $N_{1}<\dots<N_{\nu}$ for those norms, $M_{j}$ for
the sum of the $m_{i}$ at level $N_{j}$, and
$r_{K}(N)=\#\{u\in\cO_{K}:\Nm_{K/\QQ}(u)=N\}$.
\begin{enumerate}[label=\textup{(\roman*)},leftmargin=2.6em]
\item If $\nu=n+1$ then, with $P=\prod_{k}N_{k}$, $S=\ch_{0}(E)$ and
  $D_{j}=N_{j}\prod_{k\ne j}(N_{j}-N_{k})$,
  \begin{equation}\label{eq:levelsums}
    M_{j}=(-1)^{n}\,\frac{S\,P}{D_{j}},
  \end{equation}
  and the least $S>0$ making every $M_{j}$ an integer is
  $S_{\min}=\operatorname{lcm}_{j}\bigl(|D_{j}|/\gcd(|D_{j}|,P)\bigr)$.
  Multiplicity one, \Cref{lem:splitsemireg}(i), forces
  $|M_{j}|\le r_{K}(N_{j})$ for every $j$.
\item For $4\le n\le8$, every $d\in\{1,2,3,7,11,19\}$ and every choice of
  norms below the bounds recorded in item (XXII) of \Cref{sec:verification},
  no set of $n+1$ norms and no set of $n+2$ norms satisfies those constraints.
  In that range no split object of the shape of \Cref{setup:pte} is
  semiregular, and the reason is arithmetic rather than geometric: the least
  integral scale outruns the supply of elements at the levels it demands.
\item At $n=3$ the constraints of (i) can be met. With norms at most $200$
  there are exactly eleven sets of four norms that meet them, all of them at
  $d=1$ or $d=2$, and after dividing out the common factor each has one of the
  three shapes $(1,2,3,4)$, $(1,2,4,5)$ and $(1,4,9,10)$. The smallest is
  $d=2$ at $(9,18,27,36)$, where $S=1$, $|M|=(4,6,4,1)$ and $\ch_{0}(E)=1$.
\item For all eleven, the full system \eqref{eq:pteconditions} has no
  solution: no assignment of signs to the elements of $\cO_{K}$ at those
  norms, with the level sums of \eqref{eq:levelsums}, satisfies the
  off-diagonal conditions. In particular no split object of the shape of
  \Cref{setup:pte} with norms at most $200$ is semiregular at $n=3$.
\end{enumerate}
\end{theorem}

\begin{proof}
(i) The conditions $T_{r,r}=0$ for $r=1,\dots,n$ say that the linear
functional $f\mapsto\sum_{j}M_{j}f(N_{j})$ kills $x,x^{2},\dots,x^{n}$, hence
sends every polynomial $f$ of degree at most $n$ to $S\,f(0)$. Applying it to
$f(x)=\prod_{k\ne j}(x-N_{k})$, which has degree $n$, gives
\[
  M_{j}\prod_{k\ne j}(N_{j}-N_{k})=S\,(-1)^{n}\prod_{k\ne j}N_{k},
\]
and multiplying by $N_{j}$ turns this into \eqref{eq:levelsums}. The value of
$S_{\min}$ is the definition of a least common multiple, and the bound on
$|M_{j}|$ holds because each $u$ of norm $N_{j}$ contributes at most one term
$m_{i}=\pm1$ to $M_{j}$.

(ii), (iii) and (iv) are the computations of item (XXII) of
\Cref{sec:verification}, carried out in exact integer arithmetic. For (ii) and
(iii) the computation enumerates every set of $n+1$ norms represented by $\cO_{K}$
below the bound, forms $D_{j}$ and $S_{\min}$, and tests
$|M_{j}|\le r_{K}(N_{j})$; at $\nu=n+2$ the solution space of the diagonal
conditions is two dimensional and its integer points inside the same box are
swept by running two coordinates over their own ranges. For (iv) the
off-diagonal conditions are, after grouping by norm,
\[
  \sum_{j}N_{j}^{r}A_{1,j}=0\ \ (r=0,1,2),
  \qquad
  \sum_{j}N_{j}^{r}A_{2,j}=0\ \ (r=0,1),
  \qquad
  A_{k,j}=\sum_{\Nm(u)=N_{j}}m_{u}u^{k},
\]
which are additive over the levels; the four levels are therefore split into
one half that is tabulated and one that is streamed against the table, so the
enumeration is complete and not a sample. The four level sets attacked
directly are $d=2$ at $(9,18,27,36)$, $(27,54,81,108)$ and $(33,66,99,132)$
and $d=1$ at $(5,20,45,50)$; the other four of the eight are their images
under multiplication by an element of norm $2$ or $4$, which is a bijection on
the relevant norm levels because the element counts agree, so a solution for
one would give a solution for the other. The remaining three, $d=1$ at
$(13,52,117,130)$, $(17,68,153,170)$ and $(25,50,100,125)$, have levels of
$8,8,8,16$ and $12,12,12,16$ elements. For them the scale is $S=\pm S_{\min}$,
since $|S|\ge2$ violates the bound $|M_{j}|\le r_{K}(N_{j})$ and $S\mapsto-S$
is the symmetry $m\mapsto-m$. The first three levels are streamed against the
fourth, matched through a linear hash of the ten integer coordinates of the
off-diagonal conditions, and every match is checked on all ten coordinates
exactly. A solution has all ten coordinates zero and cannot be missed, so the
search is complete; it covers $4.29\times10^{14}$ sign patterns in the
largest case and finds none.
\end{proof}

\begin{remark}[What the range means]\label{rem:pterangemeaning}
The dimension that matters is $n\ge4$, since $n=2$ and $n=3$ are settled by
\cite{Mar25a,Mar25b} and \cite{FF26}. There
\Cref{thm:pterange}(ii) says that the split route closes before any geometry
is used: the level sums cannot exist, so there is nothing to test for
semiregularity. At $n=3$, where the geometry does produce an object, the level
sums do exist and the off-diagonal conditions are what remove them, which is a
different mechanism and a reminder that (ii) is a statement about a range and
not a theorem for all $n$. What both cases share is that the obstruction is
arithmetic: it lives in the norm form of $K$ and in the supply of integers of
a given norm, not in the cohomology of $A$.
\end{remark}

The range is now beside the point, because a second mechanism, this one in
the cohomology of $A$, closes the split route for every $n$ and every set of
norms at once.

\begin{corollary}[No sum of line bundles runs either argument]\label{cor:noline}
Let $A$ be of $(K,-1,n)$-Weil type with $n\ge2$, and let
$E=\bigoplus_{i}L_{i}[p_{i}]^{\oplus m_{i}}$ be a sum of shifted line bundles,
twisted by a common class or not, whose (twisted) Chern character remains of
Hodge type to first order along $\cD_{n,n}$ and has a nonzero Weil component.
Then $\sigma_{E}$ is not injective on the image of $\operatorname{ev}_{E}$:
$E$ is neither semiregular nor semiregular in the weakened sense of
\cite[Question 11.4]{Mar25b}, and it is obstructed to first order along some
direction of the Weil family. In particular every object of \Cref{setup:pte}
satisfying \eqref{eq:pteconditions}, for every $n\ge2$, every field $K$,
every rank and every set of norms, fails both criteria: the conclusion of
\Cref{thm:pterange}, that no split object of that shape is semiregular, holds
without any bound on the norms, in every dimension, and for the weakened
criterion as well.
\end{corollary}

\begin{proof}
Write $\lambda_{i}$ for the twisted first Chern classes of the summands. If
every $\lambda_{i}$ were a real multiple of $\eta$, it would be a rational
multiple, both classes being rational, and the twisted Chern character
$\sum_{i}m_{i}e^{\lambda_{i}}$ would be a polynomial in $\eta$; its component
in degree $2n$ would then lie in $\QQ\eta^{n}$, which meets the Weil line in
zero by \Cref{prop:grading}, against the hypothesis. So some $\lambda_{i}$ is
not a real multiple of $\eta$, and \Cref{thm:linebundle} applies to the
summand $L_{i}[p_{i}]$. For an object of \Cref{setup:pte} the classes are
$\lambda_{i}=\theta(u_{i})\in R_{0}$, and $R_{0}\cap\QQ\eta=0$; no twist is
needed because $c_{1}(E)=\theta(T_{1,0})=0$ by \eqref{eq:pteconditions}, as
$(1,0)$ is not an excluded pair for $n\ge2$; and \eqref{eq:pteconditions}
gives a flat Chern character with a nonzero Weil component by
\Cref{prop:chintermsofu}.
\end{proof}

\begin{remark}[What is closed and what is not]\label{rem:nolinemeaning}
\Cref{cor:noline} ends the search that \Cref{thm:pterange} began: there is no
sum of line bundles, at any member of any Weil family in any dimension, that
could run the semiregularity argument in either form, and the arithmetic of
\Cref{thm:pterange} is now a description of how the level sums fail rather
than the reason the route fails. The reason is \Cref{prop:rigidity}(i): the
only divisor class that stays of type $(1,1)$ across the whole family is the
polarisation, so every line bundle summand off that line is obstructed in
some direction of the family, and the semiregularity map, whose composite
with the obstruction is contraction with the flat class $\ch(E)$, cannot see
it. What is not closed is the object the argument needs: one with no rank one
summand off the polarisation line, of positive rank, indecomposable enough
that its Atiyah class contracted with every direction of $T$ vanishes.
Markman's reflexive sheaves at $n=3$ are of that kind, and nothing here
says that their analogues at $n\ge4$ do not exist.
\end{remark}
